\documentclass[12pt]{extarticle}
\usepackage{array}
\usepackage[T1]{fontenc}
\usepackage[Russian]{babel}
\usepackage{booktabs}
\usepackage{graphicx}
\usepackage{spreadtab}
\usepackage{numprint}
\usepackage[margin=3cm]{geometry}
\usepackage{float}
\usepackage{xfp}
\usepackage{amsmath} 
\usepackage{indentfirst}  
\begin{document}
\author{Осипов Александр Б06-604}
\title{Лабораторная работа 1.1.6. Изучение электронного осцилографа}
\maketitle
\section{Цель}

Ознакомиться с устройством и органами управления электронного или цифрового осциллографа.
Научиться измерять амплитуды и частоты произвольных сигналов, изучить основные характеристиик осциллографа.

\section{Оборудование}
Осцилограф, генераторы электрических сигналов, соединительные кабели.

\section{Ход работы}

Для расчётов, проведенных в Jupyter Notebook используется ссылка (python)

В первом эксперименте со звукого генератора (ЗГ) подавался на осциллограф сигнал определенной частоты.
По графику на экране осциллографа определялся период $T$, с сигнала, с помощью Т вычислялась частота.


Так как на ЗГ частота задавалась вручную, можно было сравнить истинную частоту (то, что показывает ЗГ) 
и частоту, расчитанную через график осциллографа. Таким образом мы определили относительную погрешность осциллографа при
измерении частоты. 

\vspace{5mm}

\textbf{\large Таблица 1 - Результаты измерения частот ЗГ с помощью осциллографа}

\begin{table}[H]
    \resizebox{\textwidth}{!}{
    \begin{spreadtab}{{tabular}{c|c|c|c|c|c|c}}
        @ $\nu$ (ЗГ), кГц & @ T, дел. & @ TIME/DIV, мс & @T, с & @ $\nu$, кГц & @ $\delta\nu$, Гц & @ $(\nu_{\text{ГЗ}} - \nu)/\nu$ \\
        1.0283 & 24 & 0.2 & b2/5*c2 & 1/d2 & (0.02/d2)*e2 & (e2 - a2)/a2 \\
        18.6*10^(-3) & 26 & 10 & b3/5*c3 & 1/d3 & (0.02/d3) *e3 & (e3 - a3)/a3 \\
        5.6744 & 17.5 & 50*10^(-3) & b4/5*c4 & 1/d4 & (0.02/d4) *e4 & (e4 - a4)/a4 \\
        184 * 10^(-3) & 27.5 & 1 & b5/5*c5 & 1/d5 & (0.02/d5) *e5 & (e5 - a5)/a5 \\
        56.082 & 18 & 5*10^(-3) & b6/5*c6 & 1/d6 & (0.02/d6) *e6 & (e6 - a6)/a6 \\
        548.74 & 18 & 0.5*10^(-3) & b7/5*c7 & 1/d7 & (0.02/d7) *e7 & (e7 - a7)/a7 \\
    \end{spreadtab}
    }
\end{table}

Максимальная относительная ошибка в измерениях $< 3.5\%$

Таким образом, можно считать, что по осцилографу можно достаточно точно определить частоту звукового сигнала.
 
\vspace{5mm}

\textbf{\large Таблица 2 - Результат измерения амплитуд ЗГ с помощью осциллографа}

\vspace{5mm}

\begin{spreadtab}{{tabular}{c|c|c|c}}
    @ & @ U, v & @ $\delta U$ & @ $\delta U /U$ \\
    @ $U_{max}$ & 20 & 5/2 & c2/b2 \\
    @ $U_{-20}$ & 2 & 0.5/2 & c3/b3 \\
    @ $U_{-40}$ & 0.2 & 0.05/2 & c4/b4 \\
    @ $U_{min}$ & 0.011 & 0.005/2 & round(c5/b5, 3) \\
\end{spreadtab}

\vspace{5mm} 

Отношения амплитуд в Дб (python):

$U_{max}/U_{-20} = 10.0$

$U_{-20}/U_{-40} = 10.0$

$U_{-40}/U_{min} = 12.6$

В случае измерения амплитуд относительные погрешности получились большими ($> 10\%$)

Можно сделать вывод, что если для измерения амплитуды неизвестного сигнала используется
осциллограф, то надо учитывать, что по оси Y он имеет высокую погрешность.
\vspace{5mm} 

\newpage
\textbf{\large Получение фигур Лиссажу}

С помощью разных комбинаций сигналов с двух источников на осциллографе были получены следующие фигуры Лиссажу:

\vspace{5mm}
Соотношение частот 1:1 : 
\vspace{5mm}

\begin{center}
  \includegraphics[scale=0.07]{assets/lab_1/Lissagu_11.jpg}
\end{center}

\vspace{5mm}
Соотношение частот 2:1 : 
\vspace{5mm}

\begin{center}
  \includegraphics[scale=0.07]{assets/lab_1/Lissagu_21.jpg}
\end{center}

\vspace{5mm}
Соотношение частот 3:1 : 
\vspace{5mm}


\begin{center}
  \includegraphics[scale=0.07]{assets/lab_1/Lissagu_31.jpg}
\end{center}

\vspace{5mm}
Соотношение частот 3:2 : 
\vspace{5mm}

\begin{center}
  \includegraphics[scale=0.07]{assets/lab_1/Lissagu_32.jpg}
\end{center}

\vspace{5mm}
Соотношение частот 4:3 : 
\vspace{5mm}

\begin{center}
  \includegraphics[scale=0.07]{assets/lab_1/Lissagu_43.jpg}
\end{center}

\vspace{5mm}
Соотношение частот 5:2 : 
\vspace{5mm}

\begin{center}
  \includegraphics[scale=0.07]{assets/lab_1/Lissagu_52.jpg}
\end{center}

Фигуры Лиссажу возникают в результате сложения двух сигналов.
Стабильные геометрические фигуры возникают только тогда, когда частоты двух сигналов соотносятся как целые числа.

По фигуре можно определить соотношение частот по формуле:

$$\nu/\nu_0 = N_h / N_v$$
где,
$\nu/\nu_0$ - соотношение частот
$N_h$ - количество пересечений горизонтальной плоскости с фигурой (без учета узлов);
$N_v$ - количество пересечений вертикальной плоскости с фигурой (без учета узлов) 

\vspace{5mm}

\textbf{\large Таблица 3 - Изучение влияния АЧХ на искажение сигнала}

$2U_0 = 6 \text{ дел.}$

\vspace{5mm}

\begin{tabular}{c|c|c}
    $U(\nu)$, дел. & $\nu$, Гц & $K = U(\nu)/\nu$ \\
    6 & 20 & 1 \\ 
    6 & 100 & 1 \\
    6 & $10^4$ & 1 \\
    5.8 & $10^5$ & 0.97 \\
    5.8 & $10^6$ & 0.97 \\
    5.4 & $10^7$ & 0.9 \\
    3.2 & $3*10^7$ & 0.53 \\
\end{tabular}

Можно заметить, что при увеличении частоты амплитуда, которую показывает осциллограф, уменьшается. 
Это связано с тем, что у осциллографа есть полоса пропускания - частота, после которой амплитуда сигнала падает примерно до 70\%.

С помощью расчетов (python) можно получить, что коэффцициент ослабления = 0.7 при частотах, максимально близких к 
макисмальной частоте осциллографа $(\approx 3*10^7 Гц)$.

\begin{center}
  \includegraphics[scale=0.7]{assets/lab_1/АЧХ.png}
\end{center}

\section{Вывод:}

Была освоена работа с осциллографом и проведены основные измеререния, которые возможно выполнить на данном устройстве.

\end{document}
